%==============================================================
%  PROJECT REPORT — Course-specific AI tutoring prototype
%  Reykjavik University
%
%  Venue-neutral scaffold. To retarget later (arXiv / workshop /
%  conference), you usually only swap the \documentclass line and
%  a few packages — the section files stay the same.
%==============================================================

\documentclass[11pt,a4paper]{article}

% ---- Core packages -----------------------------------------
\usepackage[utf8]{inputenc}
\usepackage[T1]{fontenc}
\usepackage{lmodern}
\usepackage[margin=1in]{geometry}
\usepackage{microtype}

% ---- Math & symbols ----------------------------------------
\usepackage{amsmath,amssymb,amsthm}

% ---- Figures, tables, code ---------------------------------
\usepackage{graphicx}
\usepackage{booktabs}
\usepackage{float}
\usepackage{listings}
\usepackage{xcolor}

% ---- References & links ------------------------------------
\usepackage[numbers,sort&compress]{natbib}
\usepackage[hidelinks]{hyperref}
\usepackage{cleveref}

% ---- Listing style (for code / prompt snippets) ------------
\lstset{
  basicstyle=\ttfamily\small,
  breaklines=true,
  frame=single,
  columns=fullflexible,
  showstringspaces=false
}

% ---- Title -------------------------------------------------
\title{\textbf{TODO: Report Title}\\
  \large A course-specific retrieval system that points students\\
  back to their own materials\\[0.3em]
  \normalsize \textit{(working subtitle --- replace)}}

\author{
  Helgi Skarph\'edinsson\\
  \small Reykjavik University\\
  \small \texttt{TODO: email}
  \and
  TODO: Supervising Professor\\
  \small Reykjavik University
}

\date{\today}

%==============================================================
\begin{document}

\maketitle

\begin{abstract}
\noindent
As AI tutors become more capable of producing correct answers, they risk making the students who use them less effective learners: a student who struggles towards an answer more often than not learns more from the process than one who is given the answer for free. This introduces a paradoxical problem where the goal is no longer to make the model as capable as possible, but rather to guide it toward making students effective learners. I present a course-scoped learning assistant that only answers questions within the scope of the course they are asked from, working from that course's own material such as slides and textbooks. Instead of giving answers straight back to the student, the system filters hundreds of pages of course material down to the two or three references and images that actually answer the question, pointing the student back to the exact sources they were expected to study from. In an informal trial with a small number of peer users, while they genuinely appreciated the more grounded answers within the scope of their course, and getting direct course material in responses, I still observed that they wanted to take the easy route of the direct answer they would get from other models. The results imply that while grounding answers in course material is valuable, it might not always be enough for students who are too ``addicted'' to getting answers easily. This was never about replacing ChatGPT, Gemini or Claude for students, but rather about giving a viable option to the students who do want to learn the material for themselves, rather than just have it handed to them.
\end{abstract}

% ---- Body: modular section files ---------------------------
\section{Introduction}
\label{sec:intro}

% GOAL OF THIS SECTION: motivate the problem, state the gap, state what
% you built, and preview the contribution. ~1 page.

% --- The motivating tension (you already have strong wording for this
%     in _SHARED/SOP_master_draft.md — reuse and expand it here) ---
% A student handed the answer without doing the work learns less than
% one who struggled toward it. As AI tutors get better at producing
% correct answers, they may get worse for the students using them.
TODO: open with the learning-vs-answering tension.

% --- The gap in existing tools ---
% The problem is not access to AI tools, but the absence of guidance on
% using them without letting the tool do the thinking.
TODO: state the gap precisely.

% --- Your reframing ---
% The objective shifts from "minimize loss until the model is as capable
% as possible" to "constrain and guide the model toward making students
% more effective learners."
TODO: state the reframed objective.

% --- What you built (one paragraph, high level) ---
% A course-specific tutoring prototype that works through hundreds of
% pages of course material (slides, textbooks) and filters it down to the
% two or three references and images that actually answer the student's
% question, based on what was taught in that specific course --- pointing
% students back to the exact sources they were expected to study.
TODO: one-paragraph system overview.

% --- Contributions (bullet the 3-4 concrete things this report delivers) ---
\paragraph{Contributions.}
\begin{itemize}
  \item TODO: the problem framing / design principle.
  \item TODO: the system design.
  \item TODO: the implementation.
  \item TODO: the evaluation / findings (preliminary is fine).
\end{itemize}

\section{Background and Related Work}
\label{sec:background}

% GOAL: position your work against what exists. Keep it tied to YOUR
% problem — every cited work should earn its place by contrasting with
% or supporting your approach. ~1 page.

\subsection{AI tutors and intelligent tutoring systems}
% What prior tutoring systems optimize for, and why "give a correct
% answer" is the wrong objective for learning outcomes.
TODO.

\subsection{Retrieval-augmented generation}
% RAG / retrieval over documents. You use retrieval not to feed an
% answer generator, but to surface source references. Make that
% distinction explicit — it's your novelty hook.
TODO.

\subsection{Learning science: desirable difficulties}
% Optional but strong: the pedagogy behind "struggle helps learning"
% (e.g., desirable difficulties, retrieval practice). Grounds your
% design principle in evidence rather than intuition.
TODO.

\section{System Design}
\label{sec:system}

% GOAL: the conceptual design, independent of code. A reader should
% understand HOW it works and WHY each choice was made. ~1.5-2 pages.
% Put an architecture diagram here — it does a lot of work.

\subsection{Design principle}
% The one rule the system is built around: point back to sources,
% don't substitute for thinking. State it, then show how the pipeline
% enforces it.
TODO.

\subsection{Pipeline overview}
% Walk the path: course materials in -> indexing -> student question ->
% retrieval -> filtering to 2-3 references/images -> response that cites
% the exact sources. A figure here is worth a lot.
\begin{figure}[H]
  \centering
  % \includegraphics[width=0.9\linewidth]{figures/architecture.pdf}
  \fbox{\parbox[c][4cm][c]{0.9\linewidth}{\centering TODO: architecture diagram}}
  \caption{System architecture: from course materials to source-pointing responses.}
  \label{fig:arch}
\end{figure}

\subsection{Ingesting course material}
% How slides + textbooks get parsed, chunked, and indexed. Course-scoping.
TODO.

\subsection{Retrieval and filtering}
% How you go from a query to the 2-3 most relevant references/images,
% and how you keep it grounded in what was actually taught in that course.
TODO.

\subsection{The output decision}
% The interesting part (and the thread you carry into your SOPs): deciding
% what to return — hint vs. reference vs. nothing — and where the model's
% judgment can be trusted vs. constrained.
TODO.

\section{Implementation}
\label{sec:impl}

% GOAL: the concrete build. Stack, models, data flow, key engineering
% decisions and trade-offs. Enough that someone could reproduce the shape
% of it. ~1 page.

\subsection{Stack and components}
% Languages, frameworks, embedding model, vector store, LLM(s), UI.
TODO.

\subsection{Data and indexing}
% What course(s) it was built on, volume of material, how it was processed.
TODO.

\subsection{Key engineering decisions}
% 2-3 choices worth defending, with the trade-off you accepted.
TODO.

% Example code/prompt block — delete if unused:
% \begin{lstlisting}[language=Python,caption={Retrieval + filtering step}]
% ...
% \end{lstlisting}

\section{Evaluation}
\label{sec:eval}

% GOAL: evidence that it works (or an honest account of what you could
% and couldn't test). Preliminary results are fine for a project report —
% just be clear about scope and limits. ~1 page.

\subsection{What you evaluated and how}
% Retrieval quality? Whether returned references are the "right" ones?
% User/student feedback? Define your metric(s) and setup.
TODO.

\subsection{Results}
% Tables/figures. Even a small qualitative study or a handful of worked
% examples is worth showing.
\begin{table}[H]
  \centering
  \begin{tabular}{lcc}
    \toprule
    Metric & Value & Notes \\
    \midrule
    TODO & -- & -- \\
    \bottomrule
  \end{tabular}
  \caption{TODO: results.}
  \label{tab:results}
\end{table}

\subsection{Example interactions}
% Concrete before/after: a student question -> the 2-3 sources returned.
% These are persuasive and cheap to include.
TODO.

\section{Discussion}
\label{sec:discussion}

\subsection{What worked and what it means}
% Interpret the results against the design principle from Sec. 3.
TODO.

\subsection{Limitations}
% Be honest and specific — reviewers trust honest limitations. Course
% coverage, retrieval failure modes, evaluation scope, generalization.
TODO.

\subsection{Future work}
% The handoff line for your SOPs lives here: extending the prototype into
% something testable in real classrooms, and studying the output-decision
% problem (hint vs. reference vs. withhold) under uncertainty.
TODO.

\section{Conclusion}
\label{sec:conclusion}

% ~1 short paragraph. Restate the reframed objective (guide, don't answer),
% what you built to serve it, and the one takeaway you want a reader to keep.
% Then close on the direction forward — no new information.
TODO.


% ---- Bibliography ------------------------------------------
\bibliographystyle{plainnat}
\bibliography{references}

\end{document}

% (sync test — safe to delete this line)
